\chapter{Design}
\label{chap:Design}

% REVISION NOTE (target +2 to 2.5 pages): this chapter should cash out the bridge
% idea with one concrete end-to-end example and a clearer argument for why the
% language boundary sits where it does.

This chapter turns the background of the previous chapter into concrete design decisions. It defines the specific problem within the existing RTEMS toolchain, identifies the constraints that any replacement must satisfy, and explains the bridge-based design used to introduce Haskell gradually rather than by one large swap.

\section{Problem Formulation}
\label{sec:ProblemFormulation}

The problem addressed by this dissertation is not that the existing generator failed to produce useful RTEMS tests. The problem is that one central stage of the pipeline, the Python refiner, became difficult to restructure safely as the project grew. The \texttt{testbuilder} workflow already had working Promela models, a working SPIN-based exploration stage, and working C test generation. Changing the refiner therefore meant changing the part of the system where the risk of regression was highest.

The existing Python refiner combines several kinds of work in one place. It parses annotation lines, tracks context such as whether refinement is currently inside a \texttt{STRUCT} or \texttt{SEQ}, performs lookups in a refinement dictionary, and accumulates generated code fragments. This arrangement works, but it makes it harder to reason locally about how one line of input affects the next state of the translation.

Three properties of the original implementation were especially important. First, context is stored as mutable object state. Second, the main refinement function handles many annotation forms in one place. Third, template expansion and code accumulation are interleaved with control flow. Each of these is manageable in isolation, but together they make structural changes riskier than they need to be.

\subsection{Identified Challenges}

The first challenge was behavioural preservation. The repository already contained working models and a working Python refiner. A replacement could only be justified if it produced the same observable output for the existing model set.

The second challenge was integration. The refiner is not a standalone tool. It is called from the existing Python pipeline and depends on YAML refinement dictionaries, language templates, and per-model header fragments. Any design therefore had to fit around those existing inputs rather than invent a new workflow.

The third challenge was scope control. Replacing the entire pipeline at once would have made it difficult to distinguish language-porting issues from unrelated tooling issues. A narrower design was needed so that the hardest stateful component could be rewritten first and checked continuously against the original behaviour.

\subsection{Proposed Work}

The design adopted in this dissertation is a staged port of \texttt{refine\_command}. The Haskell implementation keeps the same external role as the Python one: it receives tokenised annotation lines, updates its refinement state, and eventually renders a test body. The surrounding Python code is retained so that existing configuration files, model assets, and generation scripts can continue to drive the process.

The central design decision is therefore not only to use Haskell, but to introduce it behind a bridge that preserves the historical Python-facing interface. That keeps the migration incremental. The new implementation can be checked with the same inputs as the old one, and output differences can be detected before they spread into the rest of the workflow.

the bridge was placed at the boundary between \texttt{spin2test.py} and \texttt{refine\_command.py} for good reason. when \texttt{spin2test.py} hands work to \texttt{refine\_command.py}, most of the surrounding setup has already been done, the correct RTEMS model files have been found, the \texttt{.spn} ouput has been split into individual spin trace lines nad the YAML refinement dictionaries specific to the RTEMS model have already been loaded. This leaves \texttt{refine\_command.py} with a more singular and specialised job. it only has to recieve configuration settings at the beginning, process the spin trace lines one by one, and then return the completed test scripts. putthing the bridge here isolates the stateful, pattern matching part of the pipeline while keeping the function of the ported file more consolidated by not also having it handle other parts of the workflow.
this design decision complements the decision to use the bridge to gradually port \texttt{refine\_command.py}. this is because when comparing the output of the new file with the original, if theres a difference then i can be more sure that it has to do with the test conversion pipeline and not with some auxiliary helper function.

Two more aggressive and involved alternatives were considered for this bridge program but were ultimately not chosen.
Firstly a full rewrite at the outset would have moved the bridge location outwards, meaning that more of the system could be interfaced with using haskell, but it would also have linked  \texttt{refine\_command.py} changes with \texttt{spin2test.py}, file orchestration and template loading. as mentioned before this would have made regressions trickier to pinpoint, having a wider area to check. 

Secondly, a Foreign Function Interface (FFI) is another appraoch to the mediating between the haskell and non haskell portions of the code, this time having haskell be called as a compiled library from haskell, like the equivalent of a native library.
this way keeps the refinment inside of one process, as opposed to have the bridge process and the refinement sub process. It wasnt chosen however as it creates a more brittle boundary where things like function signatures, data layout and memory ownership have to be decided exactly. additonally it makes debugging less transparent as the haskell function is called directly by the python code. the stdin/out message pipeline contrasts this by being quite flexible  with its command response system and being extremely visible for debugging, as every message is published to console. these reasons explain why this appraoch was chosen

% SUPERVISOR NOTE (p9): this is a second good place to underline that the bridge
% strategy was your design choice. Add ~0.5 page comparing it with the two obvious
% alternatives: full rewrite first, or tighter FFI-style integration.

\section{Overview of the Design}
\label{sec:OverviewOfDesign}

% LENGTH TARGET (+1 to 1.25 pages): add a running example that starts with one
% concrete annotation line and follows it through the bridge to the emitted C fragment.
% That gives the reader a mental model before Chapter 4 gets into implementation detail.

The overall design separates the system into three logical parts: SPIN produces annotated traces from Promela models, the bridge passes those annotations and configuration data to the Haskell refiner, and the existing Python-side tooling assembles the final test files around the rendered body. This decomposition isolates the port to the refinement stage while leaving the rest of the pipeline stable.

In the repository as it stands now the direction of control remains Python centric.
\texttt{spin2test.py} reads the generated Spin trace file, keeps only lines
that start with \texttt{@@@}, splits those lines into pieces for the test code templates, loads the
refinement dictionary from YAML, and constructs the 
\texttt{command} object. That object name is preserved so the programs that call for it do not have to change, but the implementation behind it is now the
bridge in \texttt{combined\_haskell\_bridge.py}. The bridge starts the
Haskell executable in interactive mode and copies the methods that
the existing Python pipeline expects, including
\texttt{setupLanguage}, \texttt{refineSPINLine\_main}, and
\texttt{test\_body}.

A line-oriented protocol was used between Python and Haskell because it made the interface explicit and easy to inspect. The bridge sends an initial JSON configuration, then a sequence of refinement commands, and finally requests the rendered test body. The protocol is simple enough to log and replay, which proved useful during debugging.

\begin{lstlisting}[caption={Simplified bridge exchange between Python and Haskell}, label={lst:bridge-protocol}]
INIT JSON:{...}
RESPONSE:SUCCESS:INIT JSON received
COMMAND:setupLanguage
RESPONSE:SUCCESS:setupLanguage
COMMAND:refineSPINLine:["4", "CALL", "WaitForSuspend", "2", "resumeRC"]
RESPONSE:SUCCESS:refineSPINLine
COMMAND:testBody
RESPONSE:SUCCESS:testBody:BEGIN
...
RESPONSE:SUCCESS:testBody:END
\end{lstlisting}

Here is a concrete example to help  show why this appraoch works. In the
task-manager model, \texttt{spin2test.py} can encounter the spin 
line \texttt{["4", "CALL", "WaitForSuspend", "2", "resumeRC"]}
while scanning the generated trace. Python does not interpret that
line beyond tokenising it, i.e. reducing every part of the trace line into a separate piece. The token array is sent over the
bridge, and the Haskell refiner then performs the stateful work. It begins by
recording process ID 4, then applying any switch or annotation handling
required by the runtime flags, in this case its "CALL" which doesnt require any switching, then looking up \texttt{WaitForSuspend} in
the refinement dictionary, and formatting the associated template with
the arguments \texttt{2} and \texttt{resumeRC}. 

\begin{lstlisting}[language={}, caption={Refinement template for \texttt{WaitForSuspend}}, label={lst:waitforsuspend-template}]
WaitForSuspend: |
  T_log( T_NORMAL, "Waiting for Task(%d) to Suspend", {0}, {1} );
  do {{
    {1} = ( *ctx->t_isSuspend )( {0} ? taskID[{0}] : 0xffffffff );
    rtems_task_wake_after( 100 );
  }} while ({1} != RTEMS_ALREADY_SUSPENDED);
\end{lstlisting}
here is the template that is called up by searching for "WaitForSuspend" in the refinement dictionaries. the spots where theres {number} are where the tokens are slotted in

\begin{lstlisting}[language=C, caption={Generated C after template substitution}, label={lst:waitforsuspend-generated}]
T_log( T_NORMAL, "Waiting for Task(%d) to Suspend", 2, resumeRC );
do {
  resumeRC = ( *ctx->t_isSuspend )( 2 ? taskID[2] : 0xffffffff );
  rtems_task_wake_after( 100 );
} while (resumeRC != RTEMS_ALREADY_SUSPENDED);
\end{lstlisting}

and here you can see what it looks like when the values from the spin line are subsituted in.

The important design point is that Python does not need to know how that fragment is built;
it only preserves the boundary and passes the same observable inputs
that the old refiner used.

The state after that command remains inside the Haskell process. The
generated lines are accumulated in memory as part of the runtime state
rather than being written to disk immediately. Only when Python later
calls \texttt{test\_body()} does the bridge send
\texttt{COMMAND:testBody}, receive the text body between
\texttt{BEGIN} and \texttt{END} markers, and then assemble the final
output file. It does this by writing the predefined preamble and postamble text, the rendered body, and the run fragment. This separation is why the bridge can
stay narrow: Python keeps responsibility for file-system
orchestration, while Haskell owns the stateful translation from
token array stream to generated C.

JSON plus line-oriented commands were sufficient because only two
payload shapes had to cross the boundary: one configuration object at
startup and one token array per annotation line. Everything else is a
fixed command or a return of the body text. A richer RPC framework would have
added complexity without improving the one property the dissertation
needed most, namely a boundary that could be logged, replayed, and
compared directly against the historical Python behaviour.
% SUPERVISOR NOTE (p16): "Good." Keep this protocol sketch, but add a fuller
% walkthrough around it: who sends each message, what state exists before/after it,
% and why JSON plus line commands were sufficient for this project.

This design also made differential checking a key focus of the project rather than an afterthought. Because the Haskell refiner could be exercised with the same traces as the Python original, text diffs on generated tests became the main feedback loop. That in turn made it possible to move from a line-by-line translation strategy to a more Haskell-native design without losing confidence in the observable output.

Against the two main alternatives, the subprocess bridge performed
best on the criteria that mattered most for this dissertation. It preserved the
historical Python API because callers still interacted with a familiar
\texttt{command}-style object. It was easier to debug because every
request and response could be inspected as plain text. And it supported
debugging checking because the same annotation traces could be replayed
through both implementations with minimal surrounding change. Its main
cost was explicit message passing, but that cost was acceptable because
the boundary was intentionally narrow.

% LENGTH TARGET (+0.75 to 1 page): compare the adopted subprocess bridge with the
% deeper alternatives mentioned in Chapter 2 using explicit criteria: build friction,
% debuggability, preservation of the Python API, and ease of regression checking.

\section{Summary}
\label{sec:SummaryDesign}

In summary, the design of the dissertation is conservative at the system boundary and experimental inside the refiner itself. The boundary stays stable so that behaviour can be checked continuously, while the implementation behind that boundary is free to adopt a clearer Haskell structure.